\documentclass[twocolumn]{aastex631}

\usepackage{amsmath}
\usepackage{multirow}
\usepackage{natbib}
\usepackage{graphicx} 
\usepackage{aas_macros}

\begin{document}

Scalar-tensor theories of gravity, such as DHOST, often rely on protective symmetries to eliminate fatal Ostrogradsky ghost instabilities at the classical level. However, from an Effective Field Theory (EFT) perspective, these symmetries are expected to be broken by generic quantum corrections, reintroducing the ghost and creating a tension between classical stability and quantum consistency. This paper investigates the precise cost of enforcing such a symmetry in the presence of leading-order curvature-squared corrections. We analyze a DHOST-like action extended with Gauss-Bonnet and Weyl-squared terms and, after confirming that constant couplings break the symmetry, we promote these couplings to functions of the scalar field's dynamics. By demanding the symmetry be restored, we derive and solve the conditions on these functions, finding a unique, non-trivial solution that is intrinsically linked to the classical theory's structure. A subsequent Hamiltonian analysis explicitly verifies that this exceptionally symmetric action is ghost-free by construction, possessing the requisite second-class constraints. Our results quantify the severe fine-tuning required for the protective mechanism to survive, demonstrating that while possible, it is highly unnatural. This work frames the challenge for these models not as one of possibility, but of plausibility, highlighting the profound conflict between symmetry and the EFT principle of naturalness.

\end{document}
                